% !TEX program = xelatex
% ============================================================
% SLB 技术文档：6.5.4–6.5.5 功率控制与测量量
% 规范：slb-doc-generator / tech-doc-style-guide.md
% 模块类型：通用功能
% ============================================================
\documentclass[11pt,a4paper]{ctexart}

\usepackage{amsmath,amssymb,amsfonts}
\usepackage{booktabs}
\usepackage{geometry}
\usepackage{hyperref}
\usepackage{enumitem}
\usepackage{xcolor}
\usepackage{longtable}
\usepackage{array}

\geometry{left=2.5cm,right=2.5cm,top=2.5cm,bottom=2.5cm}
\hypersetup{colorlinks=true,linkcolor=blue,urlcolor=blue}

\newcolumntype{L}[1]{>{\raggedright\arraybackslash}p{#1}}
\newcommand{\dBm}{\,\mathrm{dBm}}
\newcommand{\dB}{\,\mathrm{dB}}

\title{%
  \textbf{星闪 SLB 物理层}\\[0.4em]
  \Large 6.5.4\textasciitilde 6.5.5 功率控制与测量量技术文档\\[0.3em]
  \large 功率控制 \textbullet\ 通信链路测量量 \textbullet\ 位置测量量
}
\author{依据 T/XS 10001-2025《星闪无线通信系统 接入层\\
同步低时延宽带空口 SLB 技术要求和测试方法》整理\\
（团体标准 20250326 版）}
\date{2026 年 7 月 27 日}

\begin{document}
\maketitle
\tableofcontents
\newpage

%======================================================================
\part{总览}
%======================================================================

\section{文档范围与模块划分}

本文档对应 T/XS 10001-2025 第 6.5.4\textasciitilde 6.5.5 节，覆盖物理层\textbf{共性基础算法}中的开环功率控制与通信/位置测量量。两节不规定调度策略或上报周期，只规定\textbf{怎么算} T 符号发射功率与各测量量；工程上合并为 \texttt{slb\_phy\_pwr\_meas} 库。

\begin{longtable}{L{2.4cm}L{4.0cm}L{7.5cm}}
\caption{文档范围与模块划分\label{tab:6545-scope}} \\
\toprule
\textbf{子节号} & \textbf{子节主题} & \textbf{核心输出/行为} \\
\midrule
6.5.4 & 功率控制 & $P{=}10\log_{10}(M){+}P_o{+}P_L$；多业务线性汇总；$P_{\mathrm{MAX}}$ 限幅后每子载波 $\hat{P}_k$ \\
6.5.5.1 & RSRP & 第一训练信号有效 RE 线性平均功率；多天线取 max \\
6.5.5.2 & RSSI & 整 CP-OFDM 符号总接收功率线性平均；多天线取 max \\
6.5.5.3 & RSRQ & $\mathrm{RSRP}/\mathrm{RSSI}$（同符号） \\
6.5.5.4 & SINR & $\mathrm{RSRP}/(\mathrm{RSSI}-\mathrm{RSRP})$（同符号） \\
6.5.5.5 & CBRA & $T_{\mathrm{busy}}/T_{\mathrm{measure}}$（DT 前导 / 功率门限两分支） \\
6.5.5.6 & PMS CSI & IQ[linear]+RPL 联合编码；模式 1/2 \\
6.5.5.7 & PMS 时间差 & 两/三信号 TOA/TOD $\rightarrow T$ 或 $T_1,T_2$ \\
6.5.5.8 & PMS AoA & 水平角 $[0,360)$°、垂直角 $[0,180]$° \\
\bottomrule
\end{longtable}

\section{与相关章节的关系}

\begin{longtable}{L{2.2cm}L{1.8cm}L{9.5cm}}
\caption{与相关章节的关系\label{tab:6545-rel}} \\
\toprule
\textbf{相关节号} & \textbf{关系} & \textbf{说明} \\
\midrule
6.2.2 & 上游 & 无线帧、CP-OFDM 符号时长、RE 网格；测量窗口时间粒度 \\
6.2.3 & 上游/下游 & 训练信号/SRS/PMS 为测量与功控对象；功控结果缩放 $a_{k,l}$ 幅度 \\
6.2.4/6.2.5 & 下游 & T 符号接入/ACK/数据发射前调用 6.5.4；SINR 供 MCS \\
6.5.5.1 & 内部依赖 & RSRP 为 6.5.4 路径损耗直接输入 \\
6.6 & 下游 & 消费 PMS CSI、时间差、AoA 做 RTT/定位解算 \\
XRC/系统消息 & 上游 & 下发各业务 $P_o$、\texttt{gLinkReferenceSignalPower}、测量触发 \\
第 8 章 & 下游 & 射频能力与 $P_{\mathrm{MAX}}$ 法规上限约束功控限幅 \\
\bottomrule
\end{longtable}

\section{通用功能端到端流水线}

\begin{center}
\small
\textbf{发射功控链}：第一训练符号到达
$\rightarrow$ \textbf{6.5.5.1} RSRP
$\rightarrow$ $P_L$
$\rightarrow$ \textbf{6.5.4} 单业务 $P$ / 多业务限幅
$\rightarrow$ T 符号 RE 功率映射\\[0.35em]
\textbf{链路质量链}：同符号 RSRP+RSSI
$\rightarrow$ \textbf{6.5.5.3/4} RSRQ/SINR
$\rightarrow$ MCS/调度/HARQ\\[0.35em]
\textbf{占用感知链}：符号功率 + DT 检测
$\rightarrow$ \textbf{6.5.5.5} CBRA
$\rightarrow$ LBT/直通竞争\\[0.35em]
\textbf{位置链}：PMS 接收
$\rightarrow$ \textbf{6.5.5.6\textasciitilde 8} CSI/时间差/AoA
$\rightarrow$ 6.6 测距融合
\end{center}

\newpage
%======================================================================
\part{6.5.4\textasciitilde 6.5.5 功率控制与测量量（工程解读）}
%======================================================================

\section{协议章节对应}

本节对应 T/XS 10001-2025 第 6 章第 6.5.4、6.5.5 节，涵盖：
\begin{itemize}[nosep]
  \item \textbf{6.5.4} 功率控制——T 节点 T 符号开环初始功率、多信号汇总与 $P_{\mathrm{MAX}}$ 限幅；
  \item \textbf{6.5.5.1\textasciitilde 6.5.5.5} 通信测量量——RSRP、RSSI、RSRQ、SINR、CBRA；
  \item \textbf{6.5.5.6\textasciitilde 6.5.5.8} 位置测量量——PMS CSI、时间差、AoA。
\end{itemize}

\section{功能定义}

\begin{enumerate}[nosep]
  \item \textbf{计算} 路径损耗 $P_L=\texttt{gLinkReferenceSignalPower}-\mathrm{RSRP}$；
  \item \textbf{计算} 单业务初始总功率 $P=10\log_{10}(M)+P_o+P_L$，并对同符号多业务做线性汇总与 $P_{\mathrm{MAX}}$ 限幅；
  \item \textbf{测量} RSRP/RSSI（线性平均，多天线 max），并派生 RSRQ/SINR；
  \item \textbf{测量} CBRA：按 DT 前导或功率门限累计忙时间占比；
  \item \textbf{测量} PMS CSI（IQ+RPL）、TOA/TOD 时间差与 AoA，供 6.6 使用。
\end{enumerate}

\section{模块边界（上下游及职责划分）}

\subsection{上游模块}

\begin{longtable}{L{4.2cm}L{4.5cm}L{5.5cm}}
\caption{上游模块输入\label{tab:6545-up}} \\
\toprule
\textbf{上游} & \textbf{输入} & \textbf{说明} \\
\midrule
XRC/系统消息 & 各业务 $P_o$（NominalConfig IE） & \texttt{srs/rach/ack/data-Class1/2-P0-NominalConfig} \\
G 链路配置 & \texttt{gLinkReferenceSignalPower} & 已知参考功率（dBm） \\
同步/训练接收（6.2） & 第一训练信号频域样本 & RSRP/RSSI 对象 \\
帧调度 & $M$、$N$、$K$ & 子载波数、同符号业务数、符号总子载波数 \\
射频/法规 & $P_{\mathrm{MAX}}$（$\hat{P}_{\mathrm{MAX}}$） & T 节点最大发射功率 \\
PMS 接收（6.2.3/6.6） & 测距 IQ、TOA/TOD、阵列快照 & 位置类测量 \\
\bottomrule
\end{longtable}

\subsection{下游模块}

\begin{longtable}{L{4.2cm}L{4.5cm}L{5.5cm}}
\caption{下游模块输出\label{tab:6545-down}} \\
\toprule
\textbf{下游} & \textbf{本模块输出} & \textbf{说明} \\
\midrule
T 链路发射机/功放 & 限幅后 $\hat{P}_k$ & 写入 T 符号各子载波 \\
6.5.4 内部 & RSRP & 供 $P_L$ 计算 \\
MAC/RRC & RSRP/RSSI/RSRQ/SINR/CBRA & 链路自适应、切换、LBT \\
6.6 位置流程 & PMS CSI、$T$/$T_1,T_2$、AoA & RTT/定位融合 \\
\bottomrule
\end{longtable}

\subsection{本模块不负责的内容}

\begin{longtable}{L{5.5cm}L{8.5cm}}
\caption{本模块不负责的内容\label{tab:6545-out}} \\
\toprule
\textbf{不负责内容} & \textbf{原因 / 负责节} \\
\midrule
闭环 TPC 累积修正 & 6.5.4 为开环，无 TPC 累积项 \\
$P_o$/参考功率下发 & XRC/系统消息/RRC \\
测量触发与上报周期 & MAC/RRC \\
PMS 信号生成与资源映射 & 6.2.3、6.6 \\
RTT 距离最终解算 & 6.5.5.7 仅定义时间差；解算在 6.6 \\
\bottomrule
\end{longtable}

\section{在 SLB 通信流程中的角色}

\begin{itemize}[nosep]
  \item \textbf{T 发射}：测 RSRP $\rightarrow$ 算 $P_L$/$P$ $\rightarrow$ 限幅 $\rightarrow$ RE 功率映射；
  \item \textbf{链路质量}：同符号 RSRP/RSSI $\rightarrow$ RSRQ/SINR $\rightarrow$ MCS/调度；
  \item \textbf{信道占用}：CBRA $\rightarrow$ LBT/直通竞争；
  \item \textbf{位置服务}：PMS CSI/时间差/AoA $\rightarrow$ 6.6 融合。
\end{itemize}

%----------------------------------------------------------------------
\section{关键参数与强制要求}
%----------------------------------------------------------------------

\subsection{6.5.4 功率控制}

\textbf{适用范围}：T 节点在 T 符号上发送 SRS、DMRS、PAS、接入信息、ACK、第一类/第二类数据及其配套参考信号。

\begin{equation}
P = 10\log_{10}(M) + P_o + P_L
\label{eq:power}
\end{equation}

\begin{longtable}{L{2.2cm}L{2.2cm}L{9.5cm}}
\caption{功率公式符号说明\label{tab:6545-pwr-sym}} \\
\toprule
\textbf{符号} & \textbf{单位} & \textbf{含义} \\
\midrule
$P$ & dBm & 该 T 符号、相应子载波集合上的初始总功率 \\
$M$ & 无量纲 & 本次发送占用的子载波个数（$M\geq 1$） \\
$P_o$ & dBm & 该信号/信息目标接收功率（G 配置） \\
$P_L$ & dB & 路径损耗 \\
\bottomrule
\end{longtable}

\begin{equation}
P_L = \texttt{gLinkReferenceSignalPower} - \mathrm{RSRP}
\label{eq:pl}
\end{equation}

\begin{longtable}{L{6.5cm}L{7.5cm}}
\caption{各业务 $P_o$ 配置 IE\label{tab:6545-p0}} \\
\toprule
\textbf{业务类型} & \textbf{配置参数名} \\
\midrule
信道探测信号（SRS） & \texttt{srs-P0-NominalConfig} \\
接入信息及其 DMRS、PAS & \texttt{rach-P0-NominalConfig} \\
ACK 反馈及其 DMRS & \texttt{ack-P0-NominalConfig} \\
第一类数据及其 DMRS、PAS & \texttt{data-Class1-P0-NominalConfig} \\
第二类数据及其 DMRS、PAS & \texttt{data-Class2-P0-NominalConfig} \\
\bottomrule
\end{longtable}

\textbf{多信号同 T 符号汇总与限幅}：对各业务得线性域 $\hat{P}_i$（mW），
\begin{equation}
\hat{P} = \min\left\{\hat{P}_{\mathrm{MAX}},\; \sum_{i=0}^{N-1} \hat{P}_i \right\}
\end{equation}
若 $\sum\hat{P}_i>\hat{P}_{\mathrm{MAX}}$，则 $K$ 个子载波均分：$\hat{P}_k=\hat{P}_{\mathrm{MAX}}/K$（不再保留各业务相对比例）。

\textbf{标准算例}：\texttt{gLinkReferenceSignalPower}$=-10\,\dBm$，$\mathrm{RSRP}=-70\,\dBm$ $\Rightarrow$ $P_L=60\,\dB$；ACK 的 $P_o=-100\,\dBm$，$M=12$ $\Rightarrow$ $P\approx 10.79-100+60=-29.2\,\dBm$。

\subsection{6.5.5 通信测量量}

\begin{longtable}{L{2.2cm}L{5.5cm}L{6.5cm}}
\caption{通信测量量定义摘要\label{tab:6545-comm-meas}} \\
\toprule
\textbf{量} & \textbf{定义要点} & \textbf{强制约束} \\
\midrule
RSRP & 第一训练信号所在符号、有效 RE 线性平均（W） & 多天线取 max；线性平均后再转 dBm \\
RSSI & 整符号总接收功率线性平均 & 含干扰/泄漏/噪声；多天线 max \\
RSRQ & $\mathrm{RSRP}/\mathrm{RSSI}$ & RSRP 与 RSSI \textbf{同一测量符号} \\
SINR & $\mathrm{RSRP}/(\mathrm{RSSI}-\mathrm{RSRP})$ & 同符号；分母 $\leq 0$ 判错/钳位 \\
CBRA & $T_{\mathrm{busy}}/T_{\mathrm{measure}}$ & 见门限表~\ref{tab:6545-cbra} \\
\bottomrule
\end{longtable}

\begin{longtable}{L{4.5cm}L{5.5cm}L{4cm}}
\caption{CBRA 忙判定两分支\label{tab:6545-cbra}} \\
\toprule
\textbf{场景} & \textbf{判定依据} & \textbf{门限} \\
\midrule
检测到 DT 前导 & 第一组同步序列符号功率线性均值 & $-87\,\dBm$ \\
未检测到 DT 前导 & 单 CP-OFDM 符号总接收功率线性均值 & $-62\,\dBm$ \\
\bottomrule
\end{longtable}

检测到 DT 前导且超门限时：忙时长 = 前导 + 已解析的数据帧整个持续时长（须解析帧长）。

\subsection{6.5.5 位置测量量}

\textbf{PMS CSI}：频域 $\mathbf{y}_k=\mathbf{H}_k\mathbf{x}_k+\mathbf{n}$；上报
\begin{equation}
\mathrm{IQ[\dBm]} = 20\log_{10}\!\left(\frac{|\mathrm{IQ[linear]}|}{1024}\right) + \mathrm{RPL[\dBm]}
\label{eq:iqdbm}
\end{equation}
IQ[linear]、RPL[dBm] 均为 12 bit 有符号。模式 1：全测距帧算术平均；模式 2：第 $n$ 个测距符号按跳频信道反馈。

\textbf{时间差}：两信号——第一节点 $T=t_4-t_1$，第二节点 $T=t_3-t_2$；三信号——第一节点 $T_1=t_4-t_1$、$T_2=t_5-t_4$，第二节点 $T_1=t_3-t_2$、$T_2=t_6-t_3$。

\textbf{AoA}：坐标系 $Z$=椭球法线、$X$=正南、$Y$=正东；水平角相对 $X$ 逆时针 $[0,360)$°；垂直角相对 $Z$ $[0,180]$°。

%----------------------------------------------------------------------
\section{本节核心详解}
%----------------------------------------------------------------------

\subsection{开环功控与限幅}

\textbf{功能}：按业务 $P_o$、$M$ 与实测 $P_L$ 得到发射总功率，并在多业务共符号时封顶均分。\\
\textbf{上下文}：T 符号发射前；输入 RSRP 必须来自与 \texttt{gLinkReferenceSignalPower} 对应的参考信号。

\subsection{RSRP/RSSI 与派生量}

\textbf{功能}：功率域基础量与比值质量量。\\
\textbf{上下文}：选网、功控、$P_L$、MCS；RSRQ/SINR 禁止跨符号混用样本。

\subsection{CBRA}

\textbf{功能}：信道忙时间占比。\\
\textbf{上下文}：LBT、直通竞争；DT 分支必须解析帧长，否则忙时长偏短。

\subsection{PMS 三类量}

\textbf{功能}：CSI 编码、RTT 时间差、到达角。\\
\textbf{上下文}：仅服务 6.6；本模块不输出距离米值。

%----------------------------------------------------------------------
\section{数据类型、长度与维度}
%----------------------------------------------------------------------

\begin{longtable}{L{3.2cm}L{2.2cm}L{2.2cm}L{6.5cm}}
\caption{功控接口维度\label{tab:6545-dim-pwr}} \\
\toprule
\textbf{变量} & \textbf{方向} & \textbf{长度} & \textbf{说明} \\
\midrule
\texttt{M} & 输入 & 标量 & 占用子载波数，$\geq 1$ \\
\texttt{P0\_dBm} & 输入 & 标量 & 目标接收功率 \\
\texttt{PL\_dB} & 输入 & 标量 & 路径损耗 \\
\texttt{P\_dBm} & 输出 & 标量 & 初始总功率 \\
\texttt{P\_lin[]} & 输入 & $N$ & 各业务线性功率（mW） \\
\texttt{K} & 输入 & 标量 & 符号总子载波数 \\
\texttt{Pmax\_lin} & 输入 & 标量 & $\hat{P}_{\mathrm{MAX}}$（mW） \\
\texttt{Pk\_lin[]} & 输出 & $K$ & 每子载波功率 \\
\bottomrule
\end{longtable}

\begin{longtable}{L{3.5cm}L{4.5cm}L{6cm}}
\caption{通信/位置测量接口摘要\label{tab:6545-dim-meas}} \\
\toprule
\textbf{接口} & \textbf{输入要点} & \textbf{输出} \\
\midrule
\texttt{slb\_meas\_rsrp} & 训练 RE 线性功率、天线数 & RSRP（线性或 dBm） \\
\texttt{slb\_meas\_rssi} & 符号 RE 总功率 & RSSI \\
\texttt{slb\_meas\_rsrq/sinr} & 同符号 RSRP、RSSI & 比值 \\
\texttt{slb\_meas\_cbra} & 功率序列、DT 标志、帧长 & $\mathrm{CBRA}\in[0,1]$ \\
\texttt{slb\_meas\_pms\_csi} & 频域 IQ、模式 1/2、$N_{\mathrm{sc}}$ & IQ[linear]、RPL \\
\texttt{slb\_meas\_pms\_timediff} & $t_1\ldots t_6$、信号数、角色 & $T$ 或 $T_1,T_2$ \\
\texttt{slb\_meas\_pms\_aoa} & 阵列快照 & 水平角、垂直角 \\
\bottomrule
\end{longtable}

%----------------------------------------------------------------------
\section{时序关系与触发条件}
%----------------------------------------------------------------------

\begin{longtable}{L{3.5cm}L{5.5cm}L{5.5cm}}
\caption{调用触发条件\label{tab:6545-trig}} \\
\toprule
\textbf{能力} & \textbf{进入条件} & \textbf{失败/回退} \\
\midrule
RSRP & 第一训练信号符号样本就绪 & 样本不足 $\rightarrow$ \texttt{SLB\_PHY\_ERR\_MEAS} \\
路径损耗 / 单业务 $P$ & RSRP 与 $P_o$、$M$ 已知 & $M{=}0$ 或空指针 $\rightarrow$ \texttt{ERR\_PARAM} \\
多业务限幅 & 各业务 $P_i$ 已算、$K$、$P_{\mathrm{MAX}}$ 已知 & 缓冲不足 $\rightarrow$ \texttt{ERR\_BUF} \\
RSRQ/SINR & 同符号 RSRP 与 RSSI 均已得 & 符号不一致或分母 $\leq 0$ $\rightarrow$ 错误 \\
CBRA & $T_{\mathrm{measure}}$ 内功率序列可用 & DT 分支缺帧长 $\rightarrow$ 忙时长不可信 \\
PMS 类 & 6.6 触发且时间戳/快照齐全 & 角色/信号数非法 $\rightarrow$ \texttt{ERR\_PARAM} \\
\bottomrule
\end{longtable}

\textbf{依赖链（条件推进，非状态机）}：训练符号到达 $\Rightarrow$ 可测 RSRP；调度给出 $M$/业务类型 $\Rightarrow$ 可算 $P$；射频发射前 $\Rightarrow$ 必须已完成限幅。CBRA 与 PMS 由上层独立触发，可与功控链并行。

%----------------------------------------------------------------------
\section{接口表格（明确的输入/输出）}
%----------------------------------------------------------------------

\begin{longtable}{L{3.2cm}L{2.5cm}L{1.2cm}L{1.5cm}L{5.5cm}}
\caption{\texttt{slb\_pwr\_calc\_pathloss} 接口\label{tab:6545-if-pl}} \\
\toprule
\textbf{参数} & \textbf{类型} & \textbf{必需} & \textbf{来源} & \textbf{说明} \\
\midrule
\texttt{gLinkRefPwr\_dBm} & \texttt{float} & 是 & 系统配置 & G 链路参考功率 \\
\texttt{rsrp\_dBm} & \texttt{float} & 是 & 6.5.5.1 & 测得 RSRP \\
\texttt{PL\_dB} & \texttt{float *} & 是 & 调用方 & 输出 $P_L$ \\
\bottomrule
\end{longtable}

\begin{longtable}{L{4cm}L{2cm}L{8cm}}
\caption{错误码\label{tab:6545-err}} \\
\toprule
\textbf{宏} & \textbf{值} & \textbf{含义} \\
\midrule
\texttt{SLB\_PHY\_OK} & 0 & 成功 \\
\texttt{SLB\_PHY\_ERR\_PARAM} & $-1$ & 空指针、$M{=}0$、RSRQ/SINR 分母非法 \\
\texttt{SLB\_PHY\_ERR\_BUF} & $-2$ & 输出缓冲区不足 \\
\texttt{SLB\_PHY\_ERR\_MEAS} & $-3$ & 测量样本不足或符号不一致 \\
\bottomrule
\end{longtable}

%----------------------------------------------------------------------
\section{处理步骤}
%----------------------------------------------------------------------

\subsection{端到端链路 A：测量 → 功控 → 发射}

\textbf{Step 1（训练符号 $\rightarrow$ RSRP）}
\begin{itemize}[nosep]
  \item \textbf{进入}：第一训练信号所在符号频域样本可用；
  \item \textbf{动作}：各有效 RE 线性功率求平均；多天线分别平均后取 max；
  \item \textbf{失败}：有效 RE 数为 0 $\rightarrow$ \texttt{SLB\_PHY\_ERR\_MEAS}。
\end{itemize}

\textbf{Step 2（RSRP $\rightarrow$ $P_L$）}
\begin{itemize}[nosep]
  \item \textbf{进入}：\texttt{gLinkReferenceSignalPower} 与 RSRP（dBm）已知；
  \item \textbf{动作}：$P_L=\texttt{gLinkReferenceSignalPower}-\mathrm{RSRP}$；
  \item \textbf{失败}：参考功率与测量对象不一致时上层须纠正配置，本模块只做减法。
\end{itemize}

\textbf{Step 3（业务参数 $\rightarrow$ 单业务 $P$）}
\begin{itemize}[nosep]
  \item \textbf{进入}：业务类型已映射到 $P_o$，$M\geq 1$；
  \item \textbf{动作}：$P=10\log_{10}(M)+P_o+P_L$；
  \item \textbf{失败}：$M{=}0$ $\rightarrow$ \texttt{ERR\_PARAM}。
\end{itemize}

\textbf{Step 4（同符号多业务？）}
\begin{itemize}[nosep]
  \item \textbf{进入}：若 $N{=}1$，将 $P$ 转为线性后按子载波分配并结束；若 $N{>}1$，进入限幅；
  \item \textbf{动作}：各 $P_i$ 转 $\hat{P}_i$，求和后 $\hat{P}=\min\{\hat{P}_{\mathrm{MAX}},\sum\hat{P}_i\}$；超限则 $\hat{P}_k=\hat{P}_{\mathrm{MAX}}/K$；
  \item \textbf{失败}：\texttt{Pk} 缓冲 $<K$ $\rightarrow$ \texttt{ERR\_BUF}。
\end{itemize}

\textbf{Step 5（功率 $\rightarrow$ 发射）}
\begin{itemize}[nosep]
  \item \textbf{进入}：$\hat{P}_k$ 已得；
  \item \textbf{动作}：交给 T 发射机映射 RE 幅度；本模块结束；
  \item \textbf{失败}：发射机侧问题不回灌修改功控公式。
\end{itemize}

\subsection{端到端链路 B：同符号链路质量}

\textbf{Step 1}：同符号完成 RSRP 与 RSSI（线性域）。\\
\textbf{Step 2}：若符号不一致 $\rightarrow$ \texttt{ERR\_MEAS}/\texttt{ERR\_PARAM}；否则 $\mathrm{RSRQ}=\mathrm{RSRP}/\mathrm{RSSI}$。\\
\textbf{Step 3}：若 $\mathrm{RSSI}-\mathrm{RSRP}\leq 0$ $\rightarrow$ 错误或钳位；否则 $\mathrm{SINR}=\mathrm{RSRP}/(\mathrm{RSSI}-\mathrm{RSRP})$。\\
\textbf{Step 4}：多天线分别算后取 max 上报。

\subsection{端到端链路 C：CBRA}

\textbf{Step 1}：在 $T_{\mathrm{measure}}$ 内逐符号采样功率。\\
\textbf{Step 2}：若检测到 DT 前导且同步序列功率 $>-87\,\dBm$，解析数据帧长，忙时长 += 前导+数据帧。\\
\textbf{Step 3}：否则单符号总功率 $>-62\,\dBm$ 则该符号时长计入忙。\\
\textbf{Step 4}：$\mathrm{CBRA}=T_{\mathrm{busy}}/T_{\mathrm{measure}}$。

\subsection{端到端链路 D：PMS → 6.6}

\textbf{CSI}：按模式 1/2 选样本 $\rightarrow$ 估 $\mathbf{H}_k$ $\rightarrow$ 填 IQ[linear]/RPL $\rightarrow$ 式\eqref{eq:iqdbm}。\\
\textbf{时间差}：填 $t_1\ldots t_6$ $\rightarrow$ 按信号数与节点角色选公式输出 $T$ 或 $T_1,T_2$。\\
\textbf{AoA}：阵列快照 $\rightarrow$ DOA 估计 $\rightarrow$ 映射到规定坐标系输出两角。

%----------------------------------------------------------------------
\section{协议中允许本模块改变的参数}
%----------------------------------------------------------------------

\begin{longtable}{L{4.5cm}L{4cm}L{5.5cm}}
\caption{允许配置的参数及其影响\label{tab:6545-tunable}} \\
\toprule
\textbf{参数} & \textbf{来源} & \textbf{影响} \\
\midrule
$P_o$（各 NominalConfig） & XRC/系统消息 & 目标接收功率 \\
\texttt{gLinkReferenceSignalPower} & G 配置 & $P_L$ 基准 \\
$M$, $N$, $K$ & 调度 & $P$ 与限幅分配 \\
$P_{\mathrm{MAX}}$ & 射频/法规/节电 & 限幅上限 \\
测量符号/窗口 & MAC & RSRP/RSSI/CBRA 采样 \\
PMS 模式 1/2 & 6.6 & CSI 平均或符号级反馈 \\
\bottomrule
\end{longtable}

%----------------------------------------------------------------------
\section{链路调用对照}
%----------------------------------------------------------------------

\begin{longtable}{L{4.2cm}L{2.8cm}cc}
\caption{链路与 6.5.4/6.5.5 调用对照\label{tab:6545-link}} \\
\toprule
\textbf{链路/信息块} & \textbf{标准位置} & \textbf{6.5.4} & \textbf{6.5.5} \\
\midrule
SRS（T 符号） & 6.2.3 & $P$（\texttt{srs-P0}） & --- \\
随机接入（T 符号） & 6.2.4 & $P$（\texttt{rach-P0}） & RSRP（$P_L$） \\
HARQ-ACK（T 符号） & 6.2.4.10 & $P$（\texttt{ack-P0}） & --- \\
第一类数据（T 符号） & 6.2.5.2 & $P$（\texttt{data-Class1-P0}） & SINR（MCS） \\
第二类数据（T 符号） & 6.2.5.3 & $P$（\texttt{data-Class2-P0}） & SINR \\
同步/训练接收 & 6.2.2/6.2.3 & $P_L$ 输入 & RSRP、RSSI \\
直通/LBT & --- & --- & CBRA \\
PMS 测距 & 6.2.3/6.6 & --- & CSI、时间差、AoA \\
\bottomrule
\end{longtable}

\begin{longtable}{L{2.5cm}L{2cm}L{4cm}L{5.5cm}}
\caption{测量量关系速查\label{tab:6545-meas-rel}} \\
\toprule
\textbf{测量量} & \textbf{域} & \textbf{主要依赖} & \textbf{典型用途} \\
\midrule
RSRP & 功率 & 第一训练信号 & 路径损耗、功控、选网 \\
RSSI & 功率 & 整符号 & RSRQ/SINR 分母 \\
RSRQ & 比 & RSRP/RSSI & 负载下链路质量 \\
SINR & 比 & RSRP/(RSSI-RSRP) & MCS、调度 \\
CBRA & 时间占比 & 功率门限 + DT & LBT、信道占用 \\
PMS CSI & 复信道 & 测距符号 & 测距辅助 \\
PMS 时间差 & 时间 & TOA/TOD & RTT \\
AoA & 角度 & 阵列 & 定位 \\
\bottomrule
\end{longtable}

%----------------------------------------------------------------------
\section{约束与实现要点}
%----------------------------------------------------------------------

\begin{enumerate}[nosep]
  \item \textbf{【同符号约束】}：RSRQ/SINR 的 RSRP 与 RSSI 必须同一测量符号，否则返回错误；
  \item \textbf{【线性平均】}：功率先在线性域（W/mW）平均，再转 dB/dBm；禁止对 dBm 直接算术平均；
  \item \textbf{【SINR 分母】}：$\mathrm{RSSI}-\mathrm{RSRP}\leq 0$ 时返回错误或显式钳位，禁止静默输出 NaN；
  \item \textbf{【$M\geq 1$】}：$M{=}0$ 返回 \texttt{SLB\_PHY\_ERR\_PARAM}；
  \item \textbf{【参考一致性】}：功控用 RSRP 必须与 \texttt{gLinkReferenceSignalPower} 对应同一参考信号；
  \item \textbf{【CBRA 帧长】}：DT 前导分支须解析数据帧长度，否则忙时长统计错误；
  \item \textbf{【多天线 max】}：多天线测量在本模块内统一取 max 上报，避免上层重复取 max；
  \item \textbf{【限幅均分】}：超 $P_{\mathrm{MAX}}$ 后按 $K$ 子载波均分，不再保留业务间相对功率比；
  \item \textbf{【开环边界】}：不实现 TPC 闭环累积；闭环需求由上层另行扩展且不得混入本库默认路径；
  \item \textbf{【无状态机】}：功控与测量按``进入条件 $\rightarrow$ 动作 $\rightarrow$ 失败返回''推进（训练到达、调度给出 $M$、$N{>}1$ 是否限幅、同符号标志、DT 检测结果等显式条件），不维护空闲/测量/功控/发射等 FSM 跳转表，避免异常路径锁死。
\end{enumerate}

%----------------------------------------------------------------------
\section{与标准其他章节的衔接}
%----------------------------------------------------------------------

\begin{itemize}[nosep]
  \item \textbf{6.2.2}：符号时长与 RE 网格，决定测量窗口与 $K$；
  \item \textbf{6.2.3}：训练信号/SRS/PMS 为测量与功控对象；功控结果进入 RE 幅度；
  \item \textbf{6.2.4/6.2.5}：T 符号业务发射前调用 6.5.4；SINR 服务 MCS；
  \item \textbf{6.5.5.1}：RSRP 定义是 6.5.4 $P_L$ 的直接输入；
  \item \textbf{6.6}：消费 6.5.5.6\textasciitilde 6.5.5.8 完成测距/定位解算；
  \item \textbf{XRC/系统消息}：承载 $P_o$、参考功率与测量触发；
  \item \textbf{第 8 章}：$P_{\mathrm{MAX}}$ 与射频指标约束限幅上限。
\end{itemize}

\vfill
\begin{center}
\small\textcolor{gray}{精确参数与字段定义以 T/XS 10001-2025 正式标准为准；\\
本文档提供 6.5.4\textasciitilde 6.5.5 功率控制与测量量工程解读，供 SLB 实现与联调参考。}
\end{center}

\end{document}
